% SPDX-License-Identifier: Apache-2.0
%
% Razboj, a minimal GPU in TxHDL. Built by //docs:razboj. Every listing
% is included from the tree, the picture is what Razboj drew when the
% build ran it, and the waveform is that run's.
\documentclass[journal,9pt]{IEEEtran}

% The listings are whole ranges of Rust that rustfmt sets at 80 columns,
% not the 70 the house style's 6pt is for. The frame holds 76 columns at
% 6pt, so 5.6pt, which fits 80 with about 5pt to spare (issue 1033).
\newcommand{\listingsize}{\fontsize{5.6}{6.6}\selectfont}
% SPDX-License-Identifier: Apache-2.0
%
% The house preamble, shared by every document in this directory. Loaded
% after \documentclass, before \title. Keeping it in one file is what
% stops two articles drifting apart in font, listing style or figure
% style.
% Computer Modern. IEEEtran selects Times (ptm) by default, so the face has
% to be chosen explicitly.
%
% Latin Modern is the Computer Modern variety used here, and the choice is
% forced rather than aesthetic. Under T1 encoding, plain `cmr` has no Type 1
% outlines unless cm-super is installed, and it is not in the pinned
% distribution, so pdflatex falls back to EC bitmaps and every glyph in the
% PDF becomes a Type 3 font. That was measured: the first build produced a
% document with 10 Type 3 fonts and no Type 1. Latin Modern is Computer
% Modern redrawn with a full T1 repertoire and real .pfb outlines, so the
% same design comes out scalable.
\usepackage[T1]{fontenc}
\usepackage[utf8]{inputenc}
\usepackage{lmodern}
% microtype is not in the pinned TeX distribution. emergencystretch alone
% keeps the two-column measure from producing overfull lines.
\setlength{\emergencystretch}{3em}

\usepackage{amsmath}
\usepackage{amssymb}
\usepackage{amsthm}
\usepackage{booktabs}
\usepackage{array}
% enumitem is not in the pinned TeX distribution; the list spacing below
% does what \setlist would have done.
\usepackage{float}
\usepackage{xcolor}
\usepackage{listings}
\usepackage{tikz}
\usetikzlibrary{arrows.meta,positioning,fit,backgrounds,calc}
\usepackage[hidelinks,bookmarks=true]{hyperref}

% Numbered, definition-style box for a rule the reader is meant to apply.
\theoremstyle{definition}
\newtheorem{rulebox}{Rule}
\newtheorem{example}{Example}

% Unnumbered asides.
\theoremstyle{remark}
\newtheorem*{tip}{Tip}
\newtheorem*{pitfall}{Pitfall}

% Tighter lists, in place of enumitem's \setlist.
\setlength{\topsep}{2pt}
\setlength{\itemsep}{1pt}
\setlength{\parsep}{0pt}

% Code in the running text. The typewriter face under T1 ligates `<<`
% and `>>` into guillemets, so a shift printed as \code{<<} came out as
% a French quotation mark (issue 571). microtype's \DisableLigatures is
% not in the pinned distribution, but the primitive it rests on is
% pdfTeX's own: \pdfnoligatures strips every ligature from the font it
% is given, and the current font inside \texttt is the typewriter face
% at the size in use, so each size is stripped the first time \code
% sets it. Code wants no ligature at all: two quotes are two quotes.
\newcommand{\code}[1]{\texttt{\ifdefined\pdfnoligatures\pdfnoligatures\font\fi #1}}
\newcommand{\term}[1]{\emph{#1}}

% Syntax highlighting. Four roles, distinguishable in colour and still
% distinguishable in greyscale, because an IEEE article gets printed: the
% keyword is the only bold role, the comment is the only italic one, and the
% two remaining roles differ in lightness as well as in hue.
\definecolor{kwblue}{RGB}{20,60,150}
\definecolor{cmtgreen}{RGB}{90,120,90}
\definecolor{strred}{RGB}{150,50,40}
\definecolor{numgray}{gray}{0.45}
\definecolor{framegray}{gray}{0.70}
\definecolor{bgshade}{gray}{0.97}
% A darker fill for figure cells. The listing background has to stay very
% light to keep code readable; a figure cell has to be clearly shaded or the
% distinction it draws is invisible in print.
\definecolor{cellshade}{gray}{0.90}

% The listing size is the document's choice, because it depends on the
% column: a two-column article sets `\listingsize` to 6pt and keeps its
% listings within 70 columns, or to 5.6pt when it includes source that
% rustfmt sets at 80, which is what fits a column's frame (issue 1061);
% a one-column document keeps the body size and includes source files
% at 80. `//docs:frame_test` checks every listing line against its frame. Lines are never broken; a line that does not fit
% overflows the frame, which is visible, rather than wrapping, which is
% not.
\providecommand{\listingsize}{\scriptsize}

% `'` is deliberately not a string delimiter: the language uses it for loop
% labels, as Rust does, so treating it as a quote would swallow the rest of
% the line.
\lstdefinestyle{txhdl}{
  basicstyle=\ttfamily\listingsize,
  keywordstyle=\bfseries\color{kwblue},
  commentstyle=\itshape\color{cmtgreen},
  stringstyle=\color{strred},
  numberstyle=\tiny\color{numgray},
  morekeywords={mod,use,pub,const,type,struct,enum,trait,impl,fn,async,let,mut,
    match,if,else,loop,while,for,in,break,continue,return,as,await,
    module,bus,channel,transaction,protocol,pipeline,flow,seq,config,
    port,stage,pipe,instance,connect,bind,tag,domain,blackbox,foreign,
    property,role,signal,handler,true,false,
    sequence,interface,modport,implements,package,import,var,wire,func,proc,
    case,default,next,fork,branch,join,state,fsm},
  morecomment=[l]{//},
  morestring=[b]",
  showstringspaces=false,
  columns=fullflexible,
  keepspaces=true,
  breaklines=false,
  % Framed, so a listing is bounded rather than floating in the column.
  frame=single,
  rulecolor=\color{framegray},
  framesep=3pt,
  backgroundcolor=\color{bgshade},
  xleftmargin=3pt,
  xrightmargin=3pt,
  aboveskip=6pt,
  belowskip=6pt,
  % Every listing is numbered and captioned, so it can be referred to by
  % number. `caption` without `float` numbers the listing and leaves it
  % where it was written, which is what a listing in running prose needs.
  captionpos=b,
  abovecaptionskip=2pt,
  belowcaptionskip=6pt,
}

% Figures drawn as text. Framed the same way, and with the highlighting
% switched off, because a box-and-arrow diagram is not code and half its
% words turning blue reads as an accident. Setting a style adds keys rather
% than clearing the ones already set, so the three roles are neutralised by
% hand rather than by leaving them out. Program output is printed in this
% style too, and a netlist's assignment can run past the frame, so a line
% that does not fit is broken and the rest indented; source never is.
\lstdefinestyle{ascii}{
  basicstyle=\ttfamily\listingsize,
  keywordstyle=\ttfamily,
  commentstyle=\ttfamily,
  stringstyle=\ttfamily,
  columns=fullflexible,
  keepspaces=true,
  breaklines=true,
  breakindent=2em,
  frame=single,
  rulecolor=\color{framegray},
  framesep=3pt,
  backgroundcolor=\color{bgshade},
  xleftmargin=3pt,
  xrightmargin=3pt,
  aboveskip=6pt,
  belowskip=6pt,
}
\lstset{style=txhdl}

% House TikZ styles. `shorten` lives on the arrow styles rather than on each
% arrow, so every arrow in the document gets the gap, including ones added
% later. TikZ clips a path at a node's border, so without it an arrowhead
% butts against the box with no gap at all. Keep `node distance` well above
% twice the shorten, or the arrow shrinks to nothing.
\tikzset{
  box/.style={draw,rounded corners=2pt,align=center,inner sep=4pt,
              font=\footnotesize},
  boxf/.style={box,fill=bgshade},
  lbl/.style={font=\scriptsize\itshape,align=center},
  fl/.style={-{Stealth[length=4pt]},shorten >=2.5pt,shorten <=2.5pt},
  fld/.style={fl,dashed},
}


% The contents list: the class gives a section's number 2.75em, which
% a Roman numeral past thirty overruns onto the title, so the box is
% wider here, and a subsection's entry starts where a title does.
\makeatletter
\def\l@section#1#2{\addpenalty{\@secpenalty}\addvspace{1.0em plus 1pt}%
    \@tempdima 5.75em \begingroup \parindent \z@ \rightskip \@pnumwidth%
    \parfillskip-\@pnumwidth {\bfseries\leavevmode #1}\hfil\hbox to\@pnumwidth{\hss #2}\par%
    \endgroup}
\def\l@subsection{\@dottedtocline{2}{5.75em}{5em}}
\def\l@subsubsection{\@dottedtocline{3}{10.75em}{5.5em}}
\makeatother

\usepackage{graphicx}

\InputIfFileExists{buildstamp.tex}{}{}
\providecommand{\buildstamp}{Draft build (outside the Bazel flow).}

% A range of a source file, by its begin/end markers.
\newcommand{\gpupart}[4]{%
\lstinputlisting[rangebeginprefix=//\ begin\{,rangebeginsuffix=\},%
rangeendprefix=//\ end\{,rangeendsuffix=\},includerangemarker=false,%
linerange=#2-#2,caption={#3},label={#4}]{gpu/razboj/src/#1}}
% A range of the tile crate's one file, by its markers.
\newcommand{\tilepart}[3]{%
\lstinputlisting[rangebeginprefix=//\ begin\{,rangebeginsuffix=\},%
rangeendprefix=//\ end\{,rangeendsuffix=\},includerangemarker=false,%
linerange=#1-#1,caption={#2},label={#3}]{gpu/razboj/tile/lib.rs}}

\title{Razboj: A Minimal GPU in TxHDL}

\author{Filip Filmar and Dragiša Janković%
\thanks{Both authors are independent. Filip Filmar is the
corresponding author, at f@hdlfactory.com; Dragiša Janković is
reached at linkedin.com/in/dragisa-jankovic.}%
\thanks{This is the tutorial on the first design that uses the AXI
link~\cite{axi} for something other than a test, for a reader who has
met TxHDL through the step-by-step tutorial~\cite{tutorial}. Every
listing is included from the project repository \code{HDL/txhdl} at
build time, the picture is the one Razboj drew when the build ran
it, and the waveform is that run's. The cover~\cite{cover} lists
every document.}%
\thanks{The authors used a large language model, Claude, as an
assistant in exploring the concepts and in writing and constructing
the documents and the programs; every commit of the repository says
so and records the prompts that led to it, verbatim.}%
\thanks{\buildstamp}}

\markboth{Razboj: A Minimal GPU in TxHDL}{Filmar and Janković: Razboj: A Minimal GPU in TxHDL}

\begin{document}
\maketitle

\begin{abstract}
A rasterizer consumes a display list and writes pixels into frame
memory. Razboj implements a compact graphics processing unit within
approximately 270 lines of sequential hardware description. The core
evaluates one pixel per clock cycle across bounding boxes, maintains three
edge equations, and streams covered pixels across AXI4 bursts along each
scanline. A dedicated framebuffer module terminates the peripheral end of
the bus. The display list resides within this shared memory space, allowing
the rasterizer to fetch command streams over the same bus fabric. Both units
reside in the synthesizable subset, compiling to clean netlists verified
under \code{nvc} and Verilator against cycle-accurate traces. Razboj serves
as the initial hardware master for the TxHDL AXI bus architecture.
\end{abstract}

\section{What It Draws}
\label{sec:g-what}

Figure~\ref{fig:g-scene} illustrates the reference scene generated by the
rasterizer. The hardware renders pixels into framebuffer memory over an AXI
interconnect, after which the testbench reads the image back to disk. The
scene spans 64 by 64 pixels across eight display list commands: one screen
clear, three axis-aligned rectangles, three flat-shaded triangles, and one
smoothly interpolated triangle forming the illuminated hill.

\begin{figure}[htbp]
\centering
\includegraphics[width=0.62\columnwidth]{scene.png}
\caption{Reference demonstration rendered by Razboj: 8 commands, 4096
pixels, completing in 16\,329 clock cycles.}
\label{fig:g-scene}
\end{figure}

\gpupart{scene.rs}{scene}{Display list specification from
\code{gpu/razboj/src/scene.rs}.}{lst:g-scene}

\section{The Display List}
\label{sec:g-list}

The display list defines an instruction set with four primitive types. A
clear operation specifies a solid background color. A rectangle defines an
axis-aligned bounding box. A flat triangle specifies three vertex positions,
while a shaded triangle adds per-vertex color values. These variants map
directly to a Rust \code{enum} in software drivers. Pixel colors follow an
\code{0xAARRGGBB} encoding: the alpha value occupies the most significant
byte, stored in frame memory for downstream compositors. Shaded triangles
linearly interpolate red, green, and blue components while propagating vertex
zero's alpha value across all covered pixels.

\gpupart{op.rs}{op}{The instruction set: \code{Op} as written by software,
\code{Kind} as decoded by hardware, and \code{Insn} as transmitted across
the bus.}{lst:g-op}

The hardware word \code{Insn} maintains a fixed width matching the largest
command format. Lowered hardware logic slices bus vectors at static bit
indices; variable-length formats would require runtime length decoding and
barrel shifters. A clear instruction sets bounding box, vertex coordinates,
and shading gradients to zero. A rectangle leaves vertex and shading fields
at zero, while a flat triangle clears shading fields. An opcode field
indicates valid fields for each primitive. This fixed-width format simplifies
hardware decoding at the expense of memory footprint for smaller primitives.

A host-side software assembler optimizes primitives before transmission,
executing geometry transformations once per primitive rather than per pixel.

\gpupart{op.rs}{encode}{\code{Op::encode}: software primitive assembler.
It clips bounding boxes, winds vertices, and prunes degenerate
entries.}{lst:g-encode}

An auxiliary command, \code{Op::Scissor}, defines a clipping window for
subsequent primitives. Because the rasterizer traverses only the primitive's
bounding rectangle, the assembler applies scissor tests statically by
intersecting bounding boxes at compile time. A clear command within a
scissor region transforms into an equivalent bounded rectangle. Primitives
falling entirely outside the active scissor rectangle discard immediately.
The underlying hardware rasterizer operates without runtime scissor
registers.

Hardware decoding logic inspects the primitive kind to configure three
operational paths:

\begin{lstlisting}[frame=none,numbers=none]
let boxed = Bit::from(self.kind.get() != Kind::Tri)
    & Bit::from(self.kind.get() != Kind::Shaded);
let clearing = self.skind.get() == Kind::Clear;
let wx = mux(clearing, zero16, self.sx0.get().resize::<16>());
let shaded = self.kind.get() == Kind::Shaded;
let rgb = mux(shaded, shade, self.colour.get());
\end{lstlisting}

Triangles activate edge equation evaluators, shaded triangles activate
barycentric color interpolation, and clear operations substitute full-screen
dimensions for the bounding box. Rectangles bypass edge evaluation and write
constant colors across their bounding box. Extending the architecture to
lines or textured sprites requires expanding this decode stage while
preserving the core pixel step pipeline and AXI bus writer.

\section{The List in Memory}
\label{sec:g-mem}

In system memory, the display list resides as an array of 32-bit words. The
module \code{gpu/razboj/src/dl.rs} establishes the shared binary layout
between software drivers and hardware decoders.

\gpupart{dl.rs}{format}{Display list memory layout: 16 words per instruction
with aligned fields.}{lst:g-format}

Each instruction occupies 16 words (64 bytes). The address of primitive \(i\)
evaluates as the base address plus \(i \ll 6\), substituting a bit-shift for
hardware multiplication. All instruction fields align within word boundaries,
allowing software to construct records with masks and shifts while hardware
extracts bit slices directly. A 32-bit counter word precedes the command
buffer. Software populates the command list, posts a non-zero count to
signal completion, and waits. The rasterizer polls this counter, renders the
indicated number of commands, and resets the counter word to zero upon
completion.

The rasterizer front-end polls the counter, requests 16-word commands via
16-beat AXI read bursts, latches fields sequentially, and executes primitive
rasterization. Latched fields require three setup cycles before traversal
begins, as detailed in Section~\ref{sec:g-vivado}. The front-end issues one
read burst at a time, incurring a single bus round-trip latency followed by
16 continuous beats. A 16-bit counter field supports display lists containing
up to 65\,535 entries.

\section{List After List}
\label{sec:g-lists}

Continuous animation requires processing sequential display lists frame by
frame. The rasterizer reads command counts only when its \code{ring} input
asserts. On target FPGA hardware, a doorbell register drives \code{ring}
whenever the count is non-zero, preventing idle polling traffic on the shared
bus (issue 985). In standalone simulation, \code{ring} remains asserted to
allow continuous polling. Upon completing a list, the rasterizer drains all
pending pixel write acknowledgments, clears the memory count to zero, and
awaits write confirmation before resuming counter polling. Host software
detects this transition, writes the subsequent list, and updates the
counter. The hardware \code{idle} line asserts when the count clear completes
and deasserts upon encountering a new non-zero count.

These synchronization barriers preserve memory consistency. Awaiting pixel
write responses guarantees that host processes cannot inspect partially
rendered frames upon observing an idle signal. Similarly, confirming the
count reset prevents interconnects from returning stale non-zero counts to
subsequent poll cycles, avoiding redundant rendering.

\begin{figure*}[t]
\centering
\input{razboj_lists_timing.tex}
\caption{One list ending and the next beginning. As the last
pixel's write is answered the zero goes out on \code{issue} and
\code{finished} rises; \code{idle} rises when the zero is answered.
The polls that follow, on \code{issue} and \code{rdata}, read a count of
zero while the run writes the next list in, as a program would. The
next poll reads a count of two, \code{finished} and \code{idle} fall,
and the next list's first entry arrives as one burst of sixteen beats.}
\label{fig:g-lists}
\end{figure*}

\section{The List in Tiles}
\label{sec:g-tiles}

Razboj also renders scenes in 64 by 64 pixel tiles. Tiled rasterization
allows depth testing and alpha blending to operate directly against
dedicated block RAM, eliminating repetitive read-modify-write traffic to
external DDR3 memory (issue 991). The host processor executes binning in
software without hardware modifications: sorting display lists into the
spatial tiles touched by each primitive. The standalone crate
\code{razboj\_tile} in \code{gpu/razboj/tile/lib.rs} operates without standard
libraries or dynamic memory allocation, enabling the embedded Vreteno
core~\cite{vreteno} to execute binning directly.

Each tile functions as an individual scissor boundary. Binning clips
primitive bounding boxes against intersecting tiles using standard clipping
logic. A clear command expands into bounded rectangles within covered tiles.
Triangle edge functions remain invariant across tiles because equations
evaluate against identical original vertex positions. In contrast, shaded
triangle color planes evaluate relative to the upper-left pixel of the
bounding box; clipping advances initial color accumulators to the tile's
origin using 32-bit arithmetic.

\tilepart{entry}{An entry clipped to a tile: a clear becomes a
rectangle, and a shaded triangle's planes are stepped to the tile's
first pixel.}{lst:g-tile-entry}

A rendered frame comprises a table of tile headers, each storing a
two-word descriptor followed by the sequence of intersecting primitives.

\tilepart{format}{The tile table definition from \code{razboj\_tile}.}{lst:g-tile-format}

Verification suites evaluate all test scenes across both linear and tiled
paths, confirming byte-identical results. Tests validate boundary-spanning
geometry, complex triangles intersecting multiple tiles, clipped scissor
boxes, and irregular meshes. Evaluating color planes afresh within each tile
would introduce numerical rounding divergence across seams, motivating exact
plane stepping.

The rasterizer hardware renders tiled structures directly (issue 1255).
Driver software places the tile table at the standard display list address,
offsets primitive records by two kilobytes, and asserts bit 31 of the
counter word. The rasterizer iterates through tile descriptors, rendering
primitives into a local 4096-word block RAM buffer with occupancy tracking
marks. The engine then writes the resolved tile to DDR3 memory as 64-beat
burst scanlines. Byte strobes assert only for marked pixels, leaving
untouched background locations intact. Occupancy markers store the current
tile's generation serial, avoiding costly multi-cycle buffer clears. Marker
memory scrubs every 255 tiles across 4096 clock cycles upon serial overflow.
Clearing bit 31 preserves backward compatibility by rasterizing unbinned
display lists directly to memory.

On target hardware models, rendering a 640 by 480 screen takes 326\,973
cycles in linear mode and 636\,144 cycles in single-buffered tiled mode;
rendering an icosahedron frame requires 452\,600 and 764\,607 cycles
respectively. The single-buffered prototype serializes tile rasterization
and writeback. A future double-buffered architecture will overlap rendering
with writeback while accelerating depth testing.

Hardware depth testing integrates directly into the tile buffer (issue 992).
Word 15 of an instruction enables depth testing, selects from OpenGL's eight
comparison operators, and configures depth buffer updates. Depth-enabled
primitives append a second 16-word record holding base depth values and
fractional step gradients. The rasterizer evaluates depth equations
synchronously with spatial traversal. A pipeline stage queries local depth
memory, evaluates comparisons on the following cycle, and conditionally
updates depth and color storage. Depth markers initialize to maximum distance
without full-memory clears, keeping block RAM single-ported. Linear display
lists ignore depth attributes, while tiled lists execute depth comparisons
matching architectural models bit for bit.

Blending, the alpha test and the colour mask come next
(issue~\#993). Bit 13 of word 15 says an entry has any of them, and
gives it the same second slot, whose words 3 and 4 hold the blend's two
factors, GL ES's eleven, the alpha test's comparison and reference, and
the channels written. A pixel of such an entry takes five cycles. It
reads the depth and the colour already there, the colour from a copy
of the colour bank, written where the bank is and read only by the
walk, so that each of the two keeps one write and one read and stays a
block RAM. In the next three turns it is decided, as GL orders it:
the alpha test and the depth test with the blend's eight factors; each
channel's two products and their sum; and the sum over 255, without a
divider, and the mask. A turn after that it is written. The blend had
one turn at first, and missed the board's clock by 2.2\,ns, through
twenty-two levels of logic; a pixel that tests only depth keeps its
three cycles. A pixel the mask leaves no channel of writes no colour, which
is how a clear of the depth alone is drawn. A blend over a pixel no
entry of the tile drew needs what the framebuffer holds there. So the
binning marks a tile whose entries read the colour before anything has
covered the whole tile, and the rasteriser loads such a tile from the
framebuffer first, as one more entry: a burst of the tile's width a
row, which costs about what the tile's write-out does. A frame that
begins with a clear loads nothing. Scenes under every factor are drawn
in tiles as the model draws them, byte for byte, loaded and not, and
\code{//docs:razboj\_blend} replays the netlist's load and blend under
both simulators.

Textures are next (issue~\#997), and so far they are a format and a
model, not hardware. A texture is its levels in DDR3, each texel RGBA in
32 bits as a pixel is, stored in blocks of four by four so that the
texels a cache line holds are one burst; a table of descriptors, sixteen
words each, says where each texture is and how to filter and wrap it.
Bit 14 of word 15 says an entry is textured, and gives it the second
slot and two more. They hold $uq$, $vq$ and $q$ as planes of 64 bits,
$q$ being one over the clip $w$, so that a pixel's texel is $uq/q$
there, correct in perspective; the texture's index, its environment
and that environment's colour; and the numerators of the four
derivatives GL takes the level of detail from, each of which varies
along one axis of the box only. The model divides by $q$ with a table of
1024 reciprocals and one Newton step, and takes the level of detail as
a base-2 logarithm from a table of 256. It samples under GL's two
magnification and six minification filters, repeats or clamps, and
applies the five environments of GL ES~1.1 but \code{COMBINE}, before
the alpha test. Over a floor in perspective it is within a fiftieth of a
level of the level of detail in floating point. The rasteriser reads a
textured entry's two extra slots and, until it samples, draws the entry
as if it had no texture.

\section{The Rasteriser}
\label{sec:g-raster}

Edge equations evaluate linearly across screen coordinates. For a directed
edge from \(a\) to \(b\):
\[ e(x, y) = (b_x - a_x)(y - a_y) - (b_y - a_y)(x - a_x), \]
evaluating to zero along the edge line, positive within the primitive, and
negative outside. Coordinates use fixed-point arithmetic with 1/16th pixel
precision, sampled at pixel centers. Stepping one pixel right advances the
equation by \(-16(b_y - a_y)\); stepping one row down adds \(16(b_x - a_x)\).
The hardware tracks current values, row start baselines, and step constants
for each of the three edges.

Samples falling precisely along edge boundaries adhere to top-left fill
conventions. The setup logic subtracts one unit from starting evaluations on
non-qualifying edges, ensuring shared boundary pixels rasterize exactly once
without seams or overdraw.

Shaded triangle color interpolation follows identical linear stepping rules
across independent red, green, and blue color planes. Fractional accumulators
retain 16 bits of sub-pixel precision below the active 8-bit color channels.
The host computes initial values and directional gradients, which hardware
loads without runtime division. Pixel color values clamp between 0 and 255
against fixed-point saturation boundaries.

\gpupart{raster.rs}{state}{The rasteriser's state.}{lst:g-state}

Setup calculations require six multiplications, executing in parallel across
DSP slices during the second of three preparation cycles. Per-pixel
evaluation requires only three additions and sign evaluations, plus three
additions for interpolated color channels. Arithmetic datapaths utilize
32-bit signed two's-complement logic, accommodating vertex differences
within \(\pm 1024\) pixels.

\gpupart{raster.rs}{run}{The rasteriser execution processes.}{lst:g-run}

The rasterizer partitions into two hardware processes. The first process
monitors bus completion responses, evaluates pixel coverage, and tracks
pending writes. The second process coordinates sequential execution: polling
command counters, issuing 16-beat instruction fetches, and traversing
primitive bounding boxes scanline by scanline. Hardware compilation lowers
sequential states into state machine registers and loop counters
automatically.

Traversal advances one pixel per clock cycle along scanlines, updating edge
and color accumulators concurrently. At row boundaries, accumulators reset to
baseline values plus vertical row increments during a single turnaround cycle.
Pixel generation pauses whenever the outbound AXI channel signals
backpressure, preventing buffer overflows.

\section{Writing Over AXI}
\label{sec:g-axi}

The rasterizer acts as a hardware AXI bus master. Synthesizable hardware
clients bind directly to the underlying channel ports of \code{HostClient},
eliminating software testbench abstractions such as inboxes or future queues.

Pixel sequences emit as grouped burst transactions (issue 987). Upon
encountering the initial covered pixel of a scanline, the engine opens a
write burst matching the remaining row width up to 16 beats, bounded by
4~KiB memory page frontiers. Subsequent pixels emit sequential beats with
active byte strobes on covered pixels and inactive strobes elsewhere. This
preserves in-order address and data relationships required by AXI4. Because
triangles are convex, row coverage forms continuous spans, avoiding unneeded
trailing beats. On system models, full-screen clears at 640 by 480 finish in
327\,426 cycles (approximately one cycle per pixel), improving over
single-beat writes that required 1\,229\,346 cycles.

Because write completions require only aggregate tracking and read fetches
operate strictly one burst at a time, the engine discards allocated
transaction identifiers. Inbound responses release allocation slots
immediately: write responses return identifiers upon acknowledgment, and
read bursts release tags on their terminal beat. This bounds in-flight
transactions to a maximum of four.

\begin{figure*}[t]
\centering
\input{razboj_timing.tex}
\caption{The triangle fetched and begun. The \code{word} register counts
instruction words as the front-end fetches them, receiving all 16 words on
\code{rdata} across a single burst. Setup calculations occupy three cycles,
followed by bounding box traversal where \code{hit} marks covered pixels and
\code{inflight} monitors active writes.}
\label{fig:g-wave}
\end{figure*}

\section{The Framebuffer}
\label{sec:g-fb}

The framebuffer serves as the peripheral endpoint on the AXI bus, providing
single-word pixel storage across arbitrary burst lengths. Individual beats
update only byte lanes asserted by active strobe bits, preserving surrounding
image data.

\gpupart{fb.rs}{run}{The framebuffer implementation.}{lst:g-fb}

The module stores both pixel memories and the incoming display list. Read
requests servicing command fetches return data within the acceptance cycle.
Outbound write operations buffer until data beats arrive, honoring AXI4
serialization constraints where write data channels lack explicit
identifiers.

\section{The Netlists}
\label{sec:g-netlist}

Both the rasterizer and framebuffer compile to synthesizable netlists.
Automated regression tests simulate 16 by 16 pixel reference runs against
VHDL under \code{nvc} and Verilog under Verilator at top-level module ports.

Memory arrays initialize from configuration binaries during synthesis to
ensure testbenches read valid command streams. The testbench driver
initializes channel \code{ready} signals high on reset cycle zero,
accommodating the immediate command polling sequence of the rasterizer.

\lstinputlisting[style=ascii,caption={Execution output of
\code{//gpu/razboj:trace}.},label={lst:g-out}]{out_razboj.txt}

Lowering passes escape hardware keywords to prevent toolchain syntax
collisions. Signals named \code{next} in VHDL and \code{tri} or \code{inside}
in Verilog are renamed automatically to \code{edged} and \code{hit}.
Variable naming rules also reject expressions that re-declare register names,
ensuring unambiguous netlist synthesis.

\section{Through Vivado}
\label{sec:g-vivado}

Out-of-context synthesis for the target FPGA architecture targets 32-bit
addresses, 2-bit identifiers, and a 640 by 480 display in 1024-pixel strides
(\code{//gpu/razboj:razboj\_synth}). Place-and-route reports indicate 1683
LUTs, 1323 registers, 18 DSP blocks, and zero block RAMs, representing
roughly 1\% of available device resources. Setup multiplications map
directly to DSP slices. Color shading logic adds 373 LUTs and 386 registers
without DSP utilization, as host drivers perform division ahead of time.
Alpha channels add 14 LUTs and 9 registers, while scissor testing incurs no
hardware overhead.

The core meets timing constraints at 100\,MHz with $+1.193$\,ns worst-case
setup slack. Initial unpipelined designs evaluated two 32-bit products and
subtraction in a single cycle, failing timing by $-1.660$\,ns across an
11.66\,ns critical path (issue 1034). Pipelining setup across three
cycles—coordinate differences, parallel DSP multiplication, and sum
accumulation—reduced the critical path to 8.20\,ns while adding only 24
cycles across an eight-primitive scene.

\section{What It Is Not}
\label{sec:g-not}

Razboj prioritizes minimal hardware complexity while establishing scalable
architectural patterns. The baseline design renders one pixel per cycle (three
cycles for depth-tested pixels and five for blended ones), restricts depth
testing and blending to tiled modes, limits color interpolation to screen
space without perspective correction, and gives a shaded triangle the alpha
of its first vertex. Textures are in the format and the model, not yet in
the rasteriser. The pipeline omits
geometric clipping outside bounding boxes and issues single bursts
sequentially.

Shading responsibilities divide between software and hardware: host drivers
calculate plane equations requiring division, while hardware evaluators
execute forward additions. Offloading division to hardware in future
revisions would compress 9-word shading descriptors into 3 words without
changing downstream accumulation logic (issue 989).

Coupling a CPU core to the display list interface represents the next
evolutionary step: the network-on-chip document~\cite{noc} describes the
Vreteno core~\cite{vreteno} computing icosahedron geometry and dispatching
command buffers to Razboj over a 2D mesh.

\begin{thebibliography}{9}

\bibitem{axi}
F.~Filmar and D.~Janković, ``AXI in TxHDL,'' project repository
\code{HDL/txhdl}, \code{//docs:axi}, 2026.

\bibitem{tutorial}
F.~Filmar and D.~Janković, ``TxHDL, step by step,'' project
repository \code{HDL/txhdl}, \code{//docs:tutorial}, 2026.

\bibitem{vreteno}
F.~Filmar and D.~Janković, ``Vreteno: an RV32IMAC core in TxHDL,''
project repository \code{HDL/txhdl}, \code{//docs:vreteno}, 2026.

\bibitem{noc}
F.~Filmar and D.~Janković, ``A network on chip in TxHDL,'' project
repository \code{HDL/txhdl}, \code{//docs:noc}, 2026.

\bibitem{cover}
F.~Filmar and D.~Janković, ``TxHDL documents,'' project repository
\code{HDL/txhdl}, \code{//docs:cover}, 2026.

\end{thebibliography}

\end{document}
